\documentclass[11pt,a4paper]{article}

\usepackage[utf8]{inputenc}
\usepackage[portuguese]{babel}
\usepackage{amsmath,amssymb}
\usepackage{graphicx}
\usepackage{geometry}
\geometry{a4paper, left=2.2cm, right=2.2cm, top=2.2cm, bottom=2.2cm}
\usepackage{booktabs}
\usepackage{xcolor}
\usepackage{hyperref}
\usepackage{listings}
\usepackage{tcolorbox}
\usepackage{fancyhdr}
\usepackage{enumitem}
\usepackage{float}

\hypersetup{
    colorlinks=true,
    linkcolor=blue!70!black,
    urlcolor=teal!70!black
}

\pagestyle{fancy}
\fancyhf{}
\rhead{\small TOMSA 2026 -- Guia de Execução Passo a Passo}
\lhead{\small ISEP -- Sistemas Autónomos}
\cfoot{\thepage}

\definecolor{boxbg}{rgb}{0.95,0.97,1.0}
\definecolor{boxframe}{rgb}{0.2,0.4,0.7}
\definecolor{codebg}{rgb}{0.1,0.13,0.18}
\definecolor{codetext}{rgb}{0.9,0.95,1.0}

\lstset{
    backgroundcolor=\color{codebg},
    basicstyle=\ttfamily\small\color{codetext},
    breaklines=true,
    frame=single,
    rulecolor=\color{black!30},
    showstringspaces=false
}

\title{
    \vspace{-1.0cm}
    \textbf{INSTITUTO SUPERIOR DE ENGENHARIA DO PORTO}\\[0.2cm]
    \large Tópicos de Matemática para Sistemas Autónomos (TOMSA 2026)\\[0.4cm]
    \Large \textbf{Guia Prático Passo a Passo: Onde e Como Executar o Sistema Multi-Robô}\\[0.2cm]
    \large \textit{Instruções Claras de Terminal, Dashboard Web em Plotly, Visualização QGIS e Google Earth}
}

\author{\textbf{Geovani Lopes Debastiani} (Nº 1211756)}
\date{\today}

\begin{document}
\maketitle

\begin{tcolorbox}[colback=boxbg,colframe=boxframe,title=\textbf{Resposta Rápida: Onde Rodar o Programa?}]
O programa deve ser executado no \textbf{Terminal (linha de comandos)} do seu sistema ou no terminal integrado do VS Code / Antigravity IDE, posicionado na pasta principal do projeto:
\begin{center}
\texttt{/home/geovani/antigravityProjects/ISEP/TOMSA}
\end{center}
\textbf{Se não quiser usar o terminal}, pode simplesmente clicar duas vezes no ficheiro HTML gerado:
\begin{center}
\texttt{multi\_robot\_system/output/mission\_dashboard\_bundle.html}
\end{center}
Ele abre imediatamente no Google Chrome ou Firefox com todas as 3 vistas e controlos a funcionar!
\end{tcolorbox}

\vspace{0.4cm}

% -------------------------------------------------------------
\section{Passo 1: Abrir o Terminal no Sítio Certo}
% -------------------------------------------------------------

Tem duas opções muito simples para abrir o terminal no local correto:

\subsection{Opção A: No próprio editor (Recomendado)}
\begin{enumerate}[leftmargin=1.5cm]
    \item No menu superior do seu editor (VS Code / Antigravity), clique em \textbf{Terminal} $\to$ \textbf{New Terminal} (ou prima o atalho \texttt{Ctrl + `}).
    \item O terminal abre na parte inferior do ecrã já dentro da pasta \texttt{TOMSA}.
    \item Para confirmar que está na pasta certa, digite:
\begin{lstlisting}
pwd
\end{lstlisting}
    Deve aparecer: \texttt{/home/geovani/antigravityProjects/ISEP/TOMSA}.
\end{enumerate}

\subsection{Opção B: No Terminal do Linux (Ubuntu / Debian)}
Abra a aplicação ``Terminal'' do Linux e navegue até à pasta com:
\begin{lstlisting}
cd /home/geovani/antigravityProjects/ISEP/TOMSA
\end{lstlisting}

% -------------------------------------------------------------
\section{Passo 2: Comandos para Executar o Programa}
% -------------------------------------------------------------

Copie e cole diretamente um dos seguintes comandos no terminal:

\subsection{1. Lançar o Dashboard Web Interativo (Plotly)}
Este é o comando principal que inicia o servidor local e abre o navegador automaticamente:
\begin{lstlisting}
PYTHONPATH=. .venv/bin/python3 multi_robot_system/main.py --web
\end{lstlisting}
\begin{itemize}[nosep]
    \item O navegador abrirá automaticamente em: \url{http://127.0.0.1:8050}.
    \item Para parar o servidor quando terminar, basta carregar em \texttt{Ctrl + C} no terminal.
\end{itemize}

\subsection{2. Exportar Ficheiros Geoespaciais (QGIS, Google Earth e CSV)}
Gera as trajetórias sincronizadas e os ficheiros de mapa:
\begin{lstlisting}
PYTHONPATH=. .venv/bin/python3 multi_robot_system/main.py --export
\end{lstlisting}
Ficheiros criados na pasta \texttt{multi\_robot\_system/output/}:
\begin{itemize}[nosep]
    \item \texttt{mission\_trajectories.geojson} (para abrir no QGIS).
    \item \texttt{mission\_trajectories\_3d.kml} (para abrir no Google Earth).
    \item \texttt{synchronized\_trajectories.csv} (telemetria alinhada).
\end{itemize}

\subsection{3. Correr os Testes Unitários de Matemática (Pytest)}
Executa a validação formal de todos os 24 testes matemáticos:
\begin{lstlisting}
PYTHONPATH=. .venv/bin/pytest multi_robot_system/tests/ -v
\end{lstlisting}

\subsection{4. Correr Tudo Automaticamente (--all)}
\begin{lstlisting}
PYTHONPATH=. .venv/bin/python3 multi_robot_system/main.py --all
\end{lstlisting}

% -------------------------------------------------------------
\section{Passo 3: Como Usar o Dashboard Web e as Teclas de Atalho}
% -------------------------------------------------------------

Quando a página web abrir no navegador, dispõe dos seguintes controlos interativos:

\begin{table}[H]
\centering
\begin{tabular}{cp{10.5cm}}
\toprule
\textbf{Tecla / Botão} & \textbf{O que Faz} \\
\midrule
\textbf{Tecla \texttt{1}} & Seleciona o \textbf{Robot 1} (AUV Alpha -- Vermelho). \\
\textbf{Tecla \texttt{2}} & Seleciona o \textbf{Robot 4} (AUV Bravo -- Verde). \\
\textbf{Tecla \texttt{3}} & Seleciona o \textbf{Robot 7} (AUV Charlie -- Azul). \\
\addlinespace
\textbf{Tecla \texttt{V}} & \textbf{Alterna entre as 3 Vistas}:
\begin{itemize}[nosep]
    \item \textbf{1. Global 3D}: Cenário com fundo do mar, superfície da água e rasto progressivo.
    \item \textbf{2. POV (1ª Pessoa)}: Câmara no bico do submarino com mira nos outros veículos.
    \item \textbf{3. Cockpit}: Painel de instrumentos com horizonte artificial, bússola, profundímetro, radar e alerta $<10\text{m}$.
\end{itemize} \\
\addlinespace
\textbf{Tecla \texttt{R}} & \textbf{Reinicia as trajetórias} (volta o replay imediatamente ao segundo zero). \\
\textbf{Barra de Espaço} & \textbf{Play / Pausa} da animação. \\
\addlinespace
\textbf{Botões de Velocidade} & Escolha entre \textbf{0.25x} (Lento), \textbf{1x} (Normal), \textbf{5x} (Rápido) ou \textbf{Ultra Rápido} (percorre a missão inteira em apenas 2 segundos). \\
\bottomrule
\end{tabular}
\caption{Guia de controlos rápidos da interface web.}
\end{table}

% -------------------------------------------------------------
\section{Passo 4: Como Abrir no QGIS e Google Earth}
% -------------------------------------------------------------

\subsection{Abrir no QGIS (GeoJSON)}
\begin{enumerate}[nosep]
    \item Abra o programa \textbf{QGIS}.
    \item No menu superior: clique em \textbf{Camada} $\to$ \textbf{Adicionar Camada} $\to$ \textbf{Adicionar Camada Vetorial}.
    \item Escolha o ficheiro: \texttt{multi\_robot\_system/output/mission\_trajectories.geojson}.
    \item As trajetórias aparecerão no mar da Aguçadoura. Se adicionar o mapa de satélite da Google ou OpenStreetMap no painel \textit{XYZ Tiles}, vê a costa portuguesa com os robôs no oceano!
\end{enumerate}

\subsection{Abrir no Google Earth (KML 3D)}
\begin{enumerate}[nosep]
    \item Abra o \textbf{Google Earth Pro} (ou aceda a \url{https://earth.google.com}).
    \item Vá a \textbf{Ficheiro} $\to$ \textbf{Abrir} $\to$ selecione \texttt{multi\_robot\_system/output/mission\_trajectories\_3d.kml}.
    \item O globo voa para Portugal e a barra de tempo permite ver os submarinos a submergir em 3D.
\end{enumerate}

\end{document}
